\documentclass[11pt]{article}
\usepackage[a4paper,margin=1in]{geometry}
\usepackage[T1]{fontenc}
\usepackage{lmodern}
\input{glyphtounicode}
\pdfgentounicode=1
\usepackage{amsmath,amssymb,amsthm}
\usepackage{mathtools}
\usepackage{graphicx}
\usepackage{booktabs}
\usepackage{hyperref}
\usepackage{doi}

\newtheorem{theorem}{Theorem}
\newtheorem{lemma}{Lemma}
\newtheorem{proposition}{Proposition}
\newtheorem{corollary}{Corollary}
\newtheorem{definition}{Definition}
\newtheorem{remark}{Remark}

\newcommand{\Heis}{\mathrm{Heis}_3}
\newcommand{\Zq}{\mathbb{Z}/q\mathbb{Z}}
\newcommand{\PiProj}{\Pi}
\newcommand{\Ndist}{N_{\mathrm{dist}}}
\newcommand{\Nb}{N_b}
\newcommand{\hb}{h_b}

\title{Projective Trajectory Branching and the de Sitter Limit}
\author{J\'er\^ome Beau}
\date{July 2026}

\begin{document}
\maketitle

\begin{abstract}
The de Sitter solution shows that accelerated expansion can be a property of empty geometry rather than a
consequence of matter content.
This note identifies the projective analogue of that statement inside the Heisenberg spectral admissibility
cascade.
The result does not claim that Heisenberg balls grow exponentially.
They do not: by the Bass--Guivarc'h formula their growth is polynomial of degree four.
The exponential component survives only after replacing endpoint counting by projectively distinguishable
trajectory counting.
We define a hierarchy of trajectory distinguishability notions and show that any endpoint-based definition
collapses to polynomial growth.
The canonical O12-compatible residual channel is $b$-only and therefore bounded by the non-backtracking
abelian-shadow entropy.
For a generic probe, the residual profile separates the $b$-sequences on the pre-saturation window.
Consequently, the distinguishable-history count is the Pell count
\[
\Nb(n)=2\Nb(n-1)+\Nb(n-2),
\qquad
\hb=\lim_{n\to\infty}\tfrac{1}{n}\log \Nb(n)=\log(1+\sqrt{2}).
\]
The equality is witnessed exactly on the sampled histories of a multi-prime campaign
($q\in\{29,61,101,211\}$) and the genericity condition on the probe is proved to fail only on a finite union
of proper real-algebraic hypersurfaces.
The full-history Gabor channel displays additional projective compression and is retained as the open
candidate for an evolving effective equation of state.
The cosmological conclusion is structural: accelerated expansion, if inherited by the cosmological branch,
is not driven by exponential growth of Heisenberg states but by exponential branching of projectively
distinguishable histories.
\end{abstract}

\medskip
\noindent\textbf{Keywords:} non-injective projection; Heisenberg group; spectral admissibility;
trajectory branching; projection entropy; de Sitter expansion; Pell numbers; dark energy.

\newpage
\tableofcontents

\newpage
\section{Motivation}
\label{sec:motivation}

The historical de Sitter solution~\cite{deSitter1917} establishes that accelerated expansion need not be
sourced by ordinary matter, and the Milne model~\cite{Milne1935} shows that even a linearly coasting empty
universe is not trivially excluded by expansion data.
In the projective setting considered here, the analogous question is whether accelerated expansion can arise
from the growth of admissible projection histories rather than from a material dark-energy component.

The issue is non-trivial because the spectral admissibility programme moved from LPS expander
graphs~\cite{Lubotzky1988,Beau2026a5} to the Heisenberg Cayley graph
$\Heis(\Zq)$~\cite{Beau2026a16,Beau2026a20}.
LPS balls have exponential growth, whereas Heisenberg balls have polynomial growth.
Therefore the two-regime expansion mechanism proposed in the relaxation paper~\cite{Beau2026a5} --- polynomial
shell growth followed by exponential branching of admissible trajectories --- cannot be transferred by
counting endpoints.
The only possible surviving exponential component must live in the redundancy of histories over endpoints.

This note isolates that component, proves where it can and cannot reside, pins its rate in the canonical
channel, and states precisely what is and is not established for the cosmological
branch~\cite{Beau2026d,Cosmochrony}.

\section{Endpoint Growth and the Heisenberg Obstruction}
\label{sec:endpoint}

\subsection{Polynomial endpoint growth}
\label{sec:bass-guivarch}

Let $B(n)$ denote the ball of radius $n$ in the discrete Heisenberg group with generating set
$S=\{X^{\pm1},Y^{\pm1}\}$.
By the Bass--Guivarc'h formula~\cite{Bass1972,Guivarch1973}, the homogeneous growth degree of the integer
Heisenberg group is
\[
D = 1+1+2\cdot 1 = 4 ,
\]
so $|B(n)| \asymp n^4$ before finite-$q$ saturation, and $|B(n)|$ saturates at $q^3$ on $\Heis(\Zq)$.

\begin{proposition}[Endpoint obstruction]
\label{prop:endpoint-obstruction}
Any distinguishability relation depending only on the endpoint in $\Heis(\Zq)$ has polynomial pre-saturation
growth and finite-$q$ saturation.
It cannot support an asymptotic de Sitter-like exponential rate.
\end{proposition}

\begin{proof}
Any endpoint-based invariant factors through the Heisenberg element $(a,b,\gamma)$ reached at time $n$.
The number of reachable values is bounded by the ball volume, which has polynomial Bass--Guivarc'h growth of
degree $4$ and saturates at $q^3$.
Both bounds have zero exponential rate.
\end{proof}

\subsection{Why the central fibre is not enough}
\label{sec:central-fibre}

The central coordinate $\gamma$ records the oriented area swept by the horizontal projection of the path.
This is a genuine history integral --- the level-two term of the Chen signature of the path~\cite{Chen1957}
--- but it is the only one the group retains.
A step-two nilpotent endpoint truncates the signature at level two: iterated path information of order three
or higher is erased by the group structure itself.
Counting the central register therefore does not change the growth class: the degree $4$ already decomposes
as $2$ (abelian displacement) $+\,2$ (area values $\asymp n^2$).

\begin{corollary}[Endpoint no-go]
\label{cor:endpoint-nogo}
Any invariant that is a function of $(a,b,\gamma)$ has zero exponential distinguishability rate before
finite-$q$ saturation.
\end{corollary}

\section{Trajectory Distinguishability}
\label{sec:distinguishability}

Throughout, a trajectory is a non-backtracking word $\gamma=(g_1,\dots,g_n)\in S^n$, $g_{k+1}\neq g_k^{-1}$,
identified with the path $u_0=e$, $u_k=u_{k-1}g_k$ in $\Heis(\Zq)$.
The number of such words is $4\cdot 3^{\,n-1}$, so the unconstrained history entropy rate is $\log 3$.

\begin{definition}[D1: endpoint distinguishability]
\label{def:d1}
Two trajectories are D1-distinguishable if their final Heisenberg endpoints are distinct.
\end{definition}

\begin{definition}[D2: fibre-register distinguishability]
\label{def:d2}
Two trajectories are D2-distinguishable if their final abelian displacement or their final central register
are distinct.
\end{definition}

\begin{definition}[D3: endpoint-functional distinguishability]
\label{def:d3}
Two trajectories are D3-distinguishable if some fixed function of their final Heisenberg endpoint separates
them.
\end{definition}

\begin{definition}[D4: projective trajectory distinguishability]
\label{def:d4}
Two trajectories $\gamma,\gamma'$ of length $n$ are D4-distinguishable if their projected trajectory profiles
differ,
\[
\PiProj(U_n(\gamma))\neq \PiProj(U_n(\gamma')) ,
\]
where $U_n(\gamma)$ is the sequence of projective data accumulated along the path, not only the endpoint.
\end{definition}

D1 and D2 are endpoint notions and D2 $=$ D1 on $\Heis(\Zq)$ by Section~\ref{sec:central-fibre}.
D3 is the full endpoint no-go class of Corollary~\ref{cor:endpoint-nogo}.
D4 is the only class capable of retaining exponential history branching, and the remainder of the note
determines how much of the $\log 3$ word entropy a concrete projective profile retains.

\section{The Canonical O12 Residual Channel}
\label{sec:canonical-channel}

\subsection{Setting}
\label{sec:setting}

Fix a prime $q$ and a central character $c\neq 0$.
The O12 pipeline~\cite{Beau2026a16} acts with the exact Weil representation $\rho_c$ of $\Heis(\Zq)$ on
$\mathbb{C}^q$,
\[
[\rho_c(a,b,\gamma)\varphi](x) = \psi_c\!\bigl(\gamma+b(x-a)\bigr)\,\varphi(x-a),
\qquad
\psi_c(t)=e^{2\pi i c t/q},
\]
builds fingerprint vectors over BFS shells, and accumulates their Gram--Schmidt span.
The D4 profile of a trajectory used here is the sequence of residual norms of a probe vector
$\rho_c(u_k)\,\varphi_0$ against the span frozen at the end of shell $k-1$, together with the conjugate
factor at $q-c$ following the pair-observable protocol~\cite{Beau2026a20,Beau2026a29}.
The reference state $\varphi_0$ (the probe) is either the uniform state (the O12 default) or a fixed generic
unit vector.

\subsection{The character lemma and the $b$-only theorem}
\label{sec:b-only}

\begin{lemma}[Character structure of the O12 span]
\label{lem:character}
Every Weil vector on the uniform state, and every O12 fingerprint vector, has exactly one nonzero Fourier
coefficient.
Explicitly, $\rho_c(a,b,\gamma)\,|\mathrm{unif}\rangle$ is a global phase times the additive character
$e_{cb}$, and a fingerprint built from probe points with $b$-coordinates $b_1,b_2,b_3$ and characters
$c_1,c_2,c_3$ is a global phase times $e_{c_1b_1+c_2b_2+c_3b_3}$.
Consequently the O12 span is a coordinate subspace in the Fourier basis: at the end of shell $k$ it equals
$\mathrm{span}\{e_m : m\in F_k\}$ for a deterministic frequency set $F_k\subseteq\Zq$.
\end{lemma}

\begin{proof}
On the uniform state, $[\rho_c(a,b,\gamma)\,\mathrm{unif}](x) = q^{-1/2}\,\psi_c(\gamma-ab)\,e^{2\pi i cbx/q}$,
a phase times $e_{cb}$.
A fingerprint is a pointwise product of three such vectors, hence a phase times the character at the sum
frequency.
Characters are exactly orthogonal, so the accumulated Gram--Schmidt span is the span of the characters hit,
independently of the enumeration order within shells.
\end{proof}

\begin{theorem}[$b$-only residual channel]
\label{thm:b-only}
For any probe $\varphi_0$, the residual norm of $\rho_c(a,b,\gamma)\,\varphi_0$ against the frozen O12 span
at the end of shell $k-1$ equals
\[
r_k(b)^2 \;=\; \sum_{m\notin F_{k-1}} \bigl|\widehat{\varphi_0}(m-cb)\bigr|^2 .
\]
It is invariant under translation in $a$ and under the central phase $\gamma$, and depends on the path point
only through $b$.
The mono-path profile of the canonical channel therefore factors through the $b$-shadow
$(b_1,\dots,b_n)$ of the trajectory.
\end{theorem}

\begin{proof}
In the Fourier basis, modulation by $b$ shifts frequencies by $cb$, translation by $a$ multiplies the $m$-th
coefficient by a phase, and $\gamma$ contributes a global phase:
$\widehat{\rho_c(a,b,\gamma)\varphi_0}(m) = e^{i\theta(a,\gamma,m)}\,\widehat{\varphi_0}(m-cb)$.
The residual against a coordinate subspace $\mathrm{span}\{e_m:m\in F_{k-1}\}$ is the coefficient mass on the
complement, which discards all phases.
\end{proof}

\begin{corollary}[Uniform-probe collapse; erased fibre]
\label{cor:uniform-collapse}
With the uniform probe, $r_k(b)\in\{0,1\}$ according to whether $cb\in F_{k-1}$, and the channel collapses to
a binary function of $b$.
Moreover, for any probe and any irreducible representation, the central charge acts as a global phase, so
$\gamma$ is unconditionally erased from every mono-path norm profile.
The erased fibre of the canonical D4 channel is exactly the central register together with the sign data of
the $X$-steps.
\end{corollary}

\subsection{Pell counting of $b$-shadows}
\label{sec:pell}

\begin{proposition}[Pell count]
\label{prop:pell}
Let $\Nb(n)$ be the number of $b$-increment sequences realised by non-backtracking trajectories of length
$n$.
Then $\Nb(1)=3$, $\Nb(2)=7$, and
\[
\Nb(n)=2\Nb(n-1)+\Nb(n-2) \qquad (n\ge 3),
\]
so $\Nb(n)=3,7,17,41,99,239,577,1393,3363,\dots$ (the numerators of the continued-fraction convergents of
$\sqrt2$~\cite{NSW1981}), and
\[
\hb=\lim_{n\to\infty}\tfrac1n\log\Nb(n)=\log(1+\sqrt2).
\]
\end{proposition}

\begin{proof}
The increment alphabet is $\{0,+1,-1\}$, emitted by $X^{\pm1}$, $Y$, $Y^{-1}$ respectively.
Merge $X$ and $X^{-1}$ into one class: the swap $X\leftrightarrow X^{-1}$ is a bijection of non-backtracking
words preserving $b$-emissions, so the two states generate the same shadow languages.
Non-backtracking then allows, after class $0$, all three symbols; after $+1$, only $\{0,+1\}$; after $-1$,
only $\{0,-1\}$.
The transfer matrix
\[
M=\begin{pmatrix}1&1&1\\ 1&1&0\\ 1&0&1\end{pmatrix},
\qquad
\Nb(n)=\mathbf{1}^{\!\top} M^{\,n-1}\mathbf{u},\quad \mathbf{u}=(1,1,1)^{\!\top},
\]
has characteristic polynomial $(\lambda-1)(\lambda^2-2\lambda-1)$.
Both $\Nb$ and the sequence defined by the Pell recursion with the same initial values satisfy the associated
three-term linear recursion and agree at $n=1,2,3$ ($3,7,17$ by direct computation), hence coincide.
The rate is the Perron root $1+\sqrt2$ of $\lambda^2-2\lambda-1$.
\end{proof}

\subsection{Generic separation and the branching theorem}
\label{sec:genericity}

\begin{lemma}[Genericity of profile separation]
\label{lem:genericity}
Fix a pre-saturation window $n\le n_{\max}$, meaning $F_k\neq\emptyset$ and $F_k\neq\Zq$ for all $k<n_{\max}$,
and assume $n_{\max}<q/2$ so that $b$-values do not wrap.
Then for every unit probe $\varphi_0$ outside a finite union of proper real-algebraic hypersurfaces, distinct
$b$-shadow sequences induce distinct residual profiles on the window.
\end{lemma}

\begin{proof}
Two distinct shadow sequences first differ at some step $k$, with values $b\neq b'$.
By Theorem~\ref{thm:b-only}, profile equality at step $k$ reads
$\sum_{m\in A}|\widehat{\varphi_0}(m)|^2=\sum_{m\in B}|\widehat{\varphi_0}(m)|^2$ with
$A=F_{k-1}^{\,c}-cb$ and $B=F_{k-1}^{\,c}-cb'$.
If $A=B$ then $F_{k-1}^{\,c}$ is invariant under the nonzero shift $c(b-b')$; since $q$ is prime, a
shift-invariant subset is $\emptyset$ or $\Zq$, both excluded pre-saturation.
Hence $A\neq B$ and the equality is a nonzero real quadratic form in $\widehat{\varphi_0}$ vanishing; on the
unit sphere its zero set is a proper real-algebraic hypersurface (it is violated by any $\varphi_0$ whose
Fourier mass sits on a frequency in the symmetric difference).
There are finitely many pairs of shadow sequences on a finite window, so the exceptional set is a finite
union of proper hypersurfaces.
\end{proof}

\begin{theorem}[Canonical trajectory-branching rate]
\label{thm:branching-rate}
For a generic probe in the sense of Lemma~\ref{lem:genericity}, the canonical O12-compatible D4 channel
satisfies, on the pre-saturation window,
\[
\Ndist(n)=\Nb(n),
\qquad
\lim_{n\to\infty}\tfrac1n\log\Ndist(n)=\log(1+\sqrt2).
\]
\end{theorem}

\begin{proof}
$\Ndist(n)\le\Nb(n)$ by Theorem~\ref{thm:b-only}, and $\Ndist(n)\ge\Nb(n)$ by Lemma~\ref{lem:genericity}.
The rate is Proposition~\ref{prop:pell}.
\end{proof}

\begin{remark}
\label{rem:witness}
The numerical campaign of Section~\ref{sec:numerics} verifies exact profile separation, at machine precision
and without binning, on all sampled histories at $q=211$: the number of distinct profiles equals the number
of distinct sampled $b$-sequences at every $n\le 12$ (Table~\ref{tab:pell}).
This witnesses that the concrete seeded probe is generic; the lemma guarantees that this is not an accident
of the seed.
\end{remark}

\section{Numerical Evidence}
\label{sec:numerics}

The campaign (script \texttt{trajectory\_branching.py}, simulation repository, directory
\texttt{spectral/trajectory-entropy}) samples $T$ non-backtracking trajectories on $\Heis(\Zq)$ for
$q\in\{29,61,101,211\}$ ($T=2\times10^4$ for $q\le101$, $T=10^5$ for $q=211$; three conjugate block pairs per
prime; fixed reported seeds), and evaluates three channels on one shared path ensemble:
V1, uniform probe with the O12 basis; V2, generic probe with the O12 basis (the canonical channel of
Section~\ref{sec:canonical-channel}); V3, generic probe with a Gabor basis built from the same probe.
Profiles are discretised with bin width $\epsilon\in\{0.1,0.05,0.02,0.01\}$ (with post-hoc refinement to
$10^{-3}$ at $q=211$), and the entropy rate is estimated by collision (R\'enyi-2) entropy, bracketed below by
the distinct-count of the accumulated scalar $I(n)=\sum_{k\le n}s_k$ and above by the per-step alphabet
product capped at $\log 3$.

\begin{figure}[t]
\centering
\includegraphics[width=\textwidth]{trajectory_branching_h.pdf}
\caption{Projective trajectory-branching entropy rates for $q=29,61,101,211$ (columns) and the three channels
(rows).
V1 is a collapse control.
V2 is the canonical O12-compatible $b$-shadow channel; the dashed red line is the proved ceiling
$\log(1+\sqrt2)$.
V3 is the full-history Gabor channel and exhibits additional projective compression before finite-$q$
saturation (dotted vertical lines; dashed black line $\log 3$).
Shaded bands are the lower/upper estimator brackets; thin lines show the $\epsilon$ sweep.}
\label{fig:trajectory-branching}
\end{figure}

\subsection{V1: collapse control}
\label{sec:v1}

The uniform O12 channel behaves exactly as Corollary~\ref{cor:uniform-collapse} predicts: symbols are binary,
$\epsilon$-independent, and carry at most a weak residue of the $b$-shadow when the character span is not yet
saturated (small $h\lesssim0.1$--$0.2$ at $q\ge101$, identically $0$ at $q\le61$ where all reachable
characters are spanned early).

\subsection{V2: canonical $b$-shadow channel}
\label{sec:v2}

At $q=211$ the O12 span never saturates on the window, and the refined estimate converges to the ceiling:
$h_2(n)=0.885,0.881,0.877,\dots$ for $n\ge5$ at $\epsilon=10^{-3}$, against
$\log(1+\sqrt2)=0.8814$.
Table~\ref{tab:pell} shows the exact mechanism: profiles and $b$-sequences match one-to-one on the whole
sample.
The late-window deficit of the sampled distinct counts against $\Nb(n)$ is a pure coupon-collector effect
once $\Nb(n)$ becomes comparable to $T$; the exact rate $\tfrac1n\log\Nb(n)$ approaches $\hb$ from above.

\begin{table}[t]
\centering
\begin{tabular}{rrrr}
\toprule
$n$ & $\Nb(n)$ (exact) & distinct $b$-sequences (sampled) & distinct profiles (sampled) \\
\midrule
1 & 3 & 3 & 3 \\
2 & 7 & 7 & 7 \\
3 & 17 & 17 & 17 \\
4 & 41 & 41 & 41 \\
5 & 99 & 99 & 99 \\
6 & 239 & 239 & 239 \\
7 & 577 & 577 & 577 \\
8 & 1393 & 1393 & 1393 \\
9 & 3363 & 3363 & 3363 \\
10 & 8119 & 8096 & 8096 \\
11 & 19601 & 18675 & 18675 \\
12 & 47321 & 37394 & 37394 \\
\bottomrule
\end{tabular}
\caption{Exact Pell count against sampled counts at $q=211$, $T=10^5$, canonical channel, machine precision
(no binning).
Profiles separate the sampled $b$-sequences exactly at every depth; the deficit for $n\ge10$ is sampling
incompleteness, not profile merging.}
\label{tab:pell}
\end{table}

\subsection{V3: full-history compression channel}
\label{sec:v3}

The Gabor channel replaces the O12 span by the span of $\{\rho_c(u)\varphi_0 : u\in B_j\}$, deduplicated by
$(a,b)$.
Its rank grows deterministically as the abelian disk count $2k^2+2k+1$ at every shell and for every block and
character tested, saturating at $k\simeq\sqrt{q/2}$; this matches the full-spark behaviour of generic Gabor
systems in prime dimension~\cite{Lawrence2005,Malikiosis2015}.
Because the basis is not Fourier-aligned, the residual depends on $(a_k,b_k)$, and the $(a,b)$-sequence
determines the whole word: the ceiling is the full $\log 3$.
The measured rate transiently reaches $0.90$--$0.96$ at $n=3$--$4$ ($q=211$, $\epsilon=10^{-3}$), then
decreases while the upper bracket stays near $\log 3$: profiles genuinely merge.
This is authentic projective compression of the full history, not a resolution artefact.
V3 is not used as a canonical result; it is retained as the natural proxy for full-history projective
compression and as the possible carrier of an evolving effective equation of state.

\section{Cosmological Interpretation}
\label{sec:cosmology}

\subsection{The corrected transfer from LPS to Heisenberg}
\label{sec:transfer}

The transfer from LPS to Heisenberg changes the location of exponential growth: it is absent from endpoint
volumes and present in projectively distinguishable histories.

\begin{enumerate}
\item Heisenberg endpoint volume is polynomial (Proposition~\ref{prop:endpoint-obstruction}), so endpoint
counting cannot drive acceleration.
\item Non-backtracking histories grow exponentially even on polynomial-growth groups: the redundancy
$3^n/n^4$ is structural.
\item Endpoint-based history compression collapses to the Bass--Guivarc'h no-go
(Corollary~\ref{cor:endpoint-nogo}).
\item Projective D4 distinguishability escapes the endpoint no-go, and the canonical O12-compatible channel
gives Pell growth with $\hb=\log(1+\sqrt2)$ (Theorem~\ref{thm:branching-rate}).
\item The full-history channel adds a compression profile before finite-$q$ saturation, candidate for an
evolving effective equation of state (Section~\ref{sec:v3}).
\end{enumerate}

\subsection{Relation to the de Sitter limit}
\label{sec:desitter}

The projective de Sitter reading is not that matter sources a negative-pressure substance.
It is that the admissible projection process carries a positive distinguishable-history entropy rate, so the
number of admissible reprojection histories grows exponentially independently of any material content ---
the structural analogue of expansion as a property of the dynamics itself rather than of a
source term~\cite{deSitter1917}.
In this reading, the de Sitter-like component is the asymptotic rate of projective history branching, and
the early window ($n\le2$, where profiles are still projectively degenerate) reproduces the decelerated
regime of the two-regime model of~\cite{Beau2026a5} before the branching regime takes over.

\subsection{Status of the equation of state}
\label{sec:eos}

The canonical $b$-shadow rate is rigid: it is a substrate-combinatorial constant, independent of $q$ and of
the spectral data, and therefore gives a clean positive asymptotic rate but not, by itself, an evolving
equation of state.
The evolving component, if present, must come from full-history compression, represented here by V3, whose
cosmological use remains conditional on a canonical identification of the full-history projective
observable.
Any confrontation with evolving-dark-energy indications~\cite{DESI2024} therefore belongs to the cosmological
branch~\cite{Beau2026d} and only after that identification; it is not claimed here.

\section{Neighbouring Theory}
\label{sec:neighbouring}

\subsection{Non-backtracking entropy and Ihara zeta}
The Pell count is adjacent to non-backtracking path combinatorics as organised by the Ihara zeta
function~\cite{Ihara1966}.
The relevant object here is not the growth of closed geodesics or of the group, but the growth of
non-backtracking words modulo a projection channel; the transfer-matrix mechanics are the same, the quotient
is new.

\subsection{Growth of groups}
The endpoint obstruction is a direct consequence of polynomial growth of nilpotent
groups~\cite{Bass1972,Guivarch1973}.
The result does not conflict with Bass--Guivarc'h theory; it exploits precisely the separation between
endpoint growth and path growth.

\subsection{Finite Gabor systems}
The full-history channel is a finite Weyl--Heisenberg (Gabor) system.
The deterministic rank sequence $1,5,13,25,41,61,\dots,q$ observed at every shell matches the abelian disk
count $2k^2+2k+1$, consistent with the generic linear independence of Gabor systems in prime
dimension~\cite{Lawrence2005} and its extensions~\cite{Malikiosis2015}.

\subsection{Causal sets and history counting}
The reading ``growth is carried by histories rather than states'' is adjacent to causal-set and sum-over-
histories intuitions~\cite{Bombelli1987}.
The distinction is that the present framework counts histories only through projective distinguishability
under a non-injective projection~\cite{Cosmochrony,Beau2026b}, not as independent ontological paths.

\section{Status and Open Problems}
\label{sec:status}

\begin{itemize}
\item Proved: Heisenberg endpoint growth is polynomial of degree $4$; endpoint-based distinguishability has
zero exponential rate (Proposition~\ref{prop:endpoint-obstruction}, Corollary~\ref{cor:endpoint-nogo}).
\item Proved: the O12 span is a Fourier coordinate subspace and the mono-path residual channel is $b$-only;
the central charge is unconditionally erased from mono-path norms (Lemma~\ref{lem:character},
Theorem~\ref{thm:b-only}, Corollary~\ref{cor:uniform-collapse}).
\item Proved: the $b$-shadow count is the Pell count with rate $\log(1+\sqrt2)$
(Proposition~\ref{prop:pell}), and profile separation holds for every probe outside a finite union of proper
hypersurfaces (Lemma~\ref{lem:genericity}), giving $\Ndist=\Nb$ generically
(Theorem~\ref{thm:branching-rate}).
\item Verified: exact profile separation on all sampled histories at $q=211$, machine precision, no binning
(Table~\ref{tab:pell}); ceiling convergence of the measured rate at $\epsilon\to0$.
\item Measured, open interpretation: the V3 full-history compression profile before saturation.
\item Open: canonical status of the full-history projective observable.
\item Open: extraction of an evolving $w(n)$ from projective compression.
\item Open: integration into the cosmological branch as a status paragraph and citation, after the above.
\end{itemize}

\section{Conclusion}
\label{sec:conclusion}

The Heisenberg transition does not kill the cosmological branching mechanism; it relocates it.
The endpoint space has polynomial growth, but the projectively distinguishable history space retains a
positive entropy rate.
For the canonical O12-compatible residual channel this rate is pinned exactly:
\[
\Ndist(n)=\Nb(n)=2\Nb(n-1)+\Nb(n-2),
\qquad
\hb=\log(1+\sqrt2),
\]
proved for a generic probe and witnessed exactly on the sampled multi-prime campaign.
This supplies the structural bridge from the spectral admissibility cascade to a de Sitter-like asymptotic
component independent of matter content.
The full-history compression profile remains the open candidate mechanism for an evolving effective equation
of state.

\section*{Acknowledgements}
Portions of the development of this note benefited from iterative interactions with large language models
used as analytical assistants.
The theoretical results and their interpretation remain the author's responsibility.

\clearpage
\bibliographystyle{plain}
\bibliography{cosmochrony-bibliography,external-refs}

\end{document}
